\cvsection{Expériences professionnelles}

\begin{cventries}
  \cventry
    {Stagiaire ingénieur de recherche en IA}
    {INRIA}
    {Bordeaux, France}
    {Mars 2026 -- Septembre 2026}
    {
      \begin{cvitems}
        \item {Conception et développement en autonomie d’une application fonctionnelle intégrant un agent IA multi-étapes pour aider les chercheurs à trouver des articles scientifiques.}
        \item {Développement de l’agent sans framework d’orchestration, avec l’API OpenAI, des prompts adaptés au cas d’usage et la gestion des appels d’outils, de l’état et de la mémoire.}
        \item {Développement du backend Python/FastAPI, d’API REST documentées via OpenAPI et d’un stockage MongoDB. Création d’une interface React/Next.js laissant l’utilisateur choisir les messages à conserver dans le contexte.}
        \item {Mise en œuvre d’un pipeline RAG avec chunking, embeddings, indexation dans ChromaDB et recherche sémantique.}
        \item {Automatisation de la collecte de métadonnées d’articles en ligne par tâches cron et transmission des résultats par webhooks.}
        \item {Déploiement de LLM avec vLLM sur un cluster HPC équipé de GPU NVIDIA et gestion de modèles locaux avec LM Studio sur Mac Studio.}
        \item {Évaluation de plusieurs LLM et quantifications sur HPC avec Slurm, selon la qualité des résultats, la latence et la taille du contexte.}
        \item {Évaluation d’OpenClaw dans un environnement Docker isolé pour analyser les risques liés à l’accès d’un agent aux outils et aux fichiers.}
        \item {Conteneurisation de composants avec Docker et mise en place d’une CI/CD sur GitHub.}
      \end{cvitems}
    }

  \cventry
    {Stagiaire ingénieur de recherche en apprentissage par renforcement}
    {Université Karunya}
    {Coimbatore, Inde}
    {Juin 2025 -- Septembre 2025}
    {
      \begin{cvitems}
        \item {Développement d’un modèle d’apprentissage par renforcement et de son environnement d’entraînement pour la navigation d’un robot hospitalier, puis tests sur Raspberry Pi.}
      \end{cvitems}
    }
\end{cventries}
